\section{Estado del Arte}

Como etapa inicial del presente Trabajo Profesional, se llevó adelante un proceso estructurado de generación y evaluación de ideas utilizando metodologías colaborativas de \textit{brainstorming} y \textit{brainwriting}. El objetivo de esta instancia fue identificar oportunidades tecnológicas viables que permitieran desarrollar una solución innovadora, alineada con problemáticas reales y con potencial de impacto tanto técnico como comercial.

Durante esta etapa se generó un conjunto de trece propuestas iniciales:

\begin{itemize}
    \item ScoreMaster
    \item CryptoWealth
    \item MedTrack
    \item AeroGuard
    \item SkyDictate
    \item DevQuest
    \item AutoCare
    \item SmartPantry
    \item NutriPlan
    \item ServiHogar
    \item CuentasClaras
    \item MarketScout
    \item SocioUnido
\end{itemize}

Las propuestas fueron evaluadas bajo distintos criterios relacionados con su viabilidad técnica, sostenibilidad económica, grado de innovación, potencial de mercado y motivación del equipo de desarrollo. A partir de una primera instancia de filtrado, se seleccionaron aquellas ideas con mayor potencial para realizar un análisis estratégico más profundo.

En una segunda etapa, se aplicó el modelo de análisis de las \textbf{5 C} (\textit{Compañía, Colaboradores, Clientes, Competidores y Contexto}) sobre las siguientes propuestas:

\begin{itemize}
    \item CryptoWealth
    \item MedTrack
    \item DevQuest
    \item SmartPantry
    \item CuentasClaras
\end{itemize}

\newpage

Este análisis permitió evaluar el contexto competitivo y comercial de cada solución, identificar oportunidades de diferenciación y comprender el entorno tecnológico y operativo en el que cada proyecto debería desenvolverse.

Posteriormente, las propuestas con mejores resultados fueron sometidas a un análisis adicional basado en las \textbf{4 P del marketing} (\textit{Producto, Plaza, Precio y Promoción}). Dicho estudio se realizó sobre las siguientes iniciativas:

\begin{itemize}
    \item DevQuest
    \item CuentasClaras
    \item SocioUnido
\end{itemize}

La utilización conjunta de ambos marcos conceptuales permitió comparar las distintas alternativas desde una perspectiva integral, considerando no solo aspectos técnicos, sino también factores de adopción, monetización, escalabilidad y posicionamiento estratégico.

Como resultado de este proceso iterativo de evaluación, \textbf{SocioUnido} fue seleccionada como la propuesta con mayor potencial para ser desarrollada como Trabajo Profesional. La decisión estuvo fundamentada principalmente en su capacidad para resolver problemáticas reales del ecosistema deportivo argentino, su viabilidad tecnológica, su posibilidad de escalamiento comercial y la incorporación de componentes innovadores vinculados con inteligencia artificial, automatización conversacional y sistemas seguros de validación offline.

Adicionalmente, se realizó un análisis competitivo y de estrategia de posicionamiento el cual permitió identificar que las soluciones actualmente utilizadas por gran parte de los clubes de fútbol presentan limitaciones significativas en términos de digitalización, automatización y experiencia del usuario. Muchos sistemas existentes se enfocan exclusivamente en tareas administrativas tradicionales, careciendo de herramientas predictivas, integración omnicanal y mecanismos robustos para operar en escenarios de alta concurrencia, como los accesos a estadios durante eventos deportivos.

En este contexto, SocioUnido se posiciona como una solución alineada con el estado actual de la ingeniería del software aplicada a plataformas SaaS modernas, incorporando arquitecturas escalables, automatización inteligente y modelos de interacción centrados en el usuario. El proyecto busca superar las limitaciones observadas en las soluciones tradicionales mediante una propuesta tecnológica diseñada específicamente para las necesidades operativas, económicas e institucionales de los clubes de fútbol argentinos.

Finalmente, toda la documentación complementaria correspondiente al proceso de ideación, incluyendo el detalle de las propuestas generadas y sus respectivos análisis estratégicos, será incorporada en la sección de Anexos del presente documento.
